\section{Introduction}
\label{sec:intro}

Over the last two decades, remote sensing has changed how archaeologists find and document sites. Airborne laser scanning (ALS, or LiDAR) records microtopography beneath forest canopy and has revealed entire settlement systems that ground survey had missed \citep{devereux2005woodland,chase2012geospatial,vinci2024review}. In Italy, the national ALS coverage acquired by the Ministry of the Environment has supported the (re)discovery of medieval fortified settlements \citep{lasaponara2009fullwave,masini2018canopy}, the large-scale mapping of hillforts \citep{fontana2022hillforts} and detailed analyses of castle sites in densely vegetated terrain \citep{corbo2024land,corbo2025groma}. At the same time, the free Copernicus Sentinel-2 archive, with a revisit time of a few days and 10--20\,m bands from the visible to the short-wave infrared \citep{drusch2012sentinel2}, has opened multitemporal analysis of crop and soil marks to anyone with an internet connection \citep{agapiou2014sentinel,tapete2018sentinel,abate2020sentinel}, and cloud platforms such as Google Earth Engine remove the need to download and store the imagery \citep{gorelick2017gee}.

Access to the data, however, is no longer the bottleneck; turning data into archaeological evidence is. A LiDAR-based study typically chains ground-point filtering, interpolation of a digital terrain model (DTM) and the computation and comparison of several relief visualisations, such as hillshading, sky-view factor, openness, local relief models and their combinations \citep{zaksek2011svf,doneus2013openness,hesse2010lrm,kokalj2019single}. No single visualisation is best for all features, and their choice and parameters strongly affect what an interpreter sees \citep{bennett2012comparison,stular2012visualization,kokalj2017guide}. Multispectral prospection adds its own chain: cloud masking, selection of the right acquisition dates, spectral indices, Tasselled Cap and principal component transforms, and the interpretation of marks whose visibility depends on crop type, phenology and rainfall \citep{kauth1976tasseled,agapiou2013window,masini2018cropmarks,kalayci2019contrast,abate2020sentinel}. In practice these steps are spread over GIS packages, dedicated toolboxes \citep{kokalj2016rvt,challis2011toolkit} and custom scripts. Recent work has made such workflows explicit and documented \citep{lozic2021workflow,corbo2024land}, but executing them still requires remote-sensing and programming skills that many field archaeologists, heritage officers and local research groups do not have.

Artificial intelligence (AI) is often presented as the way to close this gap. Deep convolutional networks have detected barrows, charcoal kilns and other topographic features in LiDAR derivatives \citep{verschoof2019rcnn,trier2019arran,somrak2020maya}, and machine-learning classifiers applied to multisensor, multitemporal satellite data have mapped thousands of mounds at a regional scale \citep{orengo2020mounds}. Reviews document a rapid growth of these methods in archaeology and cultural heritage \citep{davis2019obia,fiorucci2020mlch,bickler2021ml}. Three obstacles, however, keep them away from everyday practice. First, supervised models need large labelled datasets, which are expensive and exist only for a few regions \citep{kokalj2023mlready}. Second, models are usually released as research code rather than as tools that a non-specialist can install, run and question; their output is a score whose meaning and reliability are hard to judge. Third, evaluation is often carried out on curated test areas and seldom asks whether the target features can be resolved at all by the sensor, or whether known sites are a representative sample \citep{opitz2018trends,fontana2024sampling}. Several authors therefore argue for human--AI collaboration rather than full automation, with the archaeologist validating and correcting machine suggestions \citep{lambers2019citizen,doneus2022hitl,casini2023humanai}. The same concern for transparency and reproducibility runs through the debate on big data and computational methods in archaeology \citep{huggett2020bigdata,marwick2017reproducibility}.

We argue that what is missing is not another detector but a \emph{workbench}: a single environment in which an archaeologist with limited IT knowledge can bring LiDAR and multispectral evidence together and compare it with the record of known sites, in which AI assists rather than replaces interpretation, and in which every automatic result can be checked against ground truth and against the physical limits of the data. To our knowledge, no existing tool combines these three properties.

This paper presents \airch{}, an open-source, AI-assisted workbench for multi-source archaeological prospection. \airch{} runs locally in a web browser and is organised around the tasks of an archaeologist rather than around data formats. Its LiDAR module computes relief visualisations from DTM and first/last-return digital surface models \citep{kokalj2016rvt}, compares them side by side, and renders the terrain as an interactive 3D point cloud. Its multispectral module screens Sentinel-2 and Landsat archives for crop and soil marks inside a user-drawn area. Behind both modules, \airch{} embeds about 1,700 known sites from the archaeological map of the Province of Macerata \citep{perna2024camm} and the Potenza Valley Survey \citep{vermeulen2017pvs}, each linked to its primary source and with an explicit positional uncertainty; they are drawn on the maps and used as ground truth for self-evaluation rather than browsed as a separate database. A candidate dossier links each anomaly to its evidence, its area of interest and its verification status. AI enters \airch{} at three levels: physically motivated ranking of acquisition dates and anomaly maps; a benchmark that measures each detection layer against known sites; and a conversational assistant that suggests parameters, explains the layers in plain language and reminds the user of what the data can and cannot show. \todo{Align this sentence with the final implementation of the assistant and of any learned detector.}

The contributions of this paper are:
\begin{enumerate}
  \item \textbf{A unified, low-barrier workbench.} \airch{} integrates LiDAR visualisation, 3D exploration, multitemporal multispectral screening and candidate dossiers in one no-code interface, designed with and for archaeologists with limited IT background (Section~\ref{sec:design}).
  \item \textbf{Evidence-first AI assistance.} We describe a scene-ranking and anomaly-detection pipeline that encodes archaeological remote-sensing knowledge (feature-scale contrast, crop phase, antecedent drought, parcel-interior statistics) and an assistant that keeps the user in the loop, presenting scores as investigation priorities rather than probabilities (Section~\ref{sec:ai}).
  \item \textbf{A built-in, reproducible evaluation protocol.} \airch{} measures its own layers against known sites with AUC, precision@$k$ against a chance baseline and a field-boundary contamination index, and we report a resolvability analysis that states which classes of buried features 10\,m imagery can and cannot detect (Section~\ref{sec:evaluation}).
  \item \textbf{A user-centred evaluation.} \todo{We assess usability and usefulness with N archaeologists through task-based sessions and the System Usability Scale \citep{brooke1996sus}.}
\end{enumerate}

Applied to the Province of Macerata, the benchmark shows that a conventional anomaly pipeline is dominated by agrarian geometry: field boundaries covering 7.5\% of the analysable area attracted 70\% of the strongest candidates, and computing statistics inside parcels reduced this contamination tenfold. It also shows, with honest numbers, that the Sentinel-2 branch cannot detect site-scale features: the 5\,m streets of the Roman town of \textit{Potentia} reach a single-date signal-to-noise ratio of only 1.6, and site-level AUC stays between 0.34 and 0.58. Rather than hiding these results, \airch{} surfaces them to the user and routes site-scale questions to LiDAR and very-high-resolution imagery, reserving multispectral screening for features wider than about 20\,m, such as roads, centuriation, large enclosures and palaeochannels. We consider this transparency a design requirement for AI tools aimed at non-specialists.

The remainder of the paper is organised as follows. Section~\ref{sec:related} reviews related work on LiDAR visualisation, multispectral prospection and AI in archaeology. Section~\ref{sec:design} derives the requirements and design principles of \airch{}, and Section~\ref{sec:architecture} describes its architecture. Section~\ref{sec:ai} details the AI components. Section~\ref{sec:evaluation} presents the study area, the detection benchmark, the resolvability analysis, a LiDAR case study and the user study. Section~\ref{sec:discussion} discusses limitations and the ethics of publishing site locations, and Section~\ref{sec:conclusions} concludes.
